\subsubsection{HTML/SVG presentation of compact actual bindings}

The same readability problem that motivated compact actual-binding rendering in TXT
also arises in graphical exports. In particular, when large structured actual values are
rendered directly inside DOT/SVG node labels, the resulting graphics become difficult
to read and may expand horizontally in ways that obscure the assurance-case structure.

Accordingly, graphical export should follow the same general design principle already
adopted for TXT export:

\begin{itemize}
\item keep the main assurance-case presentation readable and structurally focused;
\item present compact actual-binding identifiers inline where needed; and
\item make the full bound values available separately in a clearly delimited details view.
\end{itemize}

For HTML export, the preferred presentation is:

\begin{itemize}
\item the main SVG or graph panel shows only compact IDs inline (for example, \#S1, \#N2),
not full concrete terms;
\item a side panel, expandable details view, or equivalent auxiliary pane presents the full
actual-value legend;
\item selecting or hovering over a compact ID in the surrounding HTML presentation may
highlight the corresponding full value in the legend, but such interaction is a presentation
enhancement rather than a semantic requirement.
\end{itemize}

For plain SVG output viewed outside an HTML container, the preferred presentation is:

\begin{itemize}
\item keep the node labels visually clean by showing only compact IDs inline;
\item do not attempt to embed the full actual legend into the visible node labels; and
\item if needed, emit the full legend separately as accompanying TXT/HTML material rather
than burdening the SVG itself.
\end{itemize}

This approach has several advantages.

\begin{itemize}
\item It preserves the readability of the argument graph by preventing large concrete terms
from dominating the node text.
\item It keeps the graphical presentation aligned with the already adopted compact-ID model
of TXT export.
\item It preserves inspectability, because the full actual values remain available in a separate,
explicit legend.
\item It avoids overfitting the DOT/SVG layer to the pretty-printing of Prolog terms, which is
not the primary purpose of the graphical view.
\end{itemize}

Operationally, this means that the graphical exporter should treat compact actual-binding
IDs as the default inline representation for structured bound values, and should regard
full-value presentation as a separate details-layer concern.

\paragraph{Status.}
This is a design decision and implementation direction. The current implementation
already supports compact actual-binding rendering and a delimited legend in TXT export.
A corresponding HTML/SVG details presentation remains future work.
